\chapter{系统设计与实现}
\label{ch:system}

本章围绕系统的实际实现细节展开，与第二章的理论描述相比，讨论的是当前实现真实采用的状态组织、操作路径与落地的替代方式。

\section{总体架构与数据路径}
系统所创建的哈希表从逻辑上看是长度为 \(N\) 的开放寻址表，而且 \(N\) 永远是 2 的幂次，其主要包含五种操作，分别是初始化，插入，查询，删除以及渐进迁移。在初始化阶段要生成 \(N\) 和 \(K\)，并形成 Facility 结构；插入阶段首先由 Router 减小候选范围，然后在 tail Cubby 中执行有界 k-kick；查询阶段借助“Router 候选 + 商压缩恢复校验”来执行命中判断；删除阶段则要释放槽位并执行必要的回填，这些操作过程会记录 router\_probe\_steps\footnote{单次操作中 Router 前缀树遍历的步数。}、kick\_count\footnote{单次插入或删除中 k-kick 链发生的槽位搬移次数。}和 cubby\_tier\footnote{本次操作所涉 Cubby 的 tiered 层级编号。}，从而可以在后续的实验统计使用。

每张表用 TableState\footnote{哈希表运行时状态，保存元素个数、表规模 \(N\)、Router 桶数 \(K\) 及 Facility 数组等。}来表示，它包含表规模 \(N\)，局部桶数 \(K\) 以及 Facility 数组，其中 Facility 负责管理多 tiered cubbies，一个唯一的 tail Cubby 和长度为 \(K\) 的 Router 数组，而 Cubby 则是实际的存储容器，Slot 是 Cubby 内的数组槽位。

查询的主路径大致如下，如算法~\ref{alg:facility-cubby-slot} 所示：首先计算 \(g_x=\pi(x)\)，按照 \(g_x\) 的低 \(\log N\) 位找到 Facility 和 Router ；Router给出候选\((\text{Cubby}^*,\textit{slot})\) 之后，接着读取槽内的 payload 和 MiniArray 元数据，重新创建 \(\pi(x^{'})\) ，并与查询键的置换值做对比。

\begin{algorithm}[htbp]
  \caption{查询过程}
  \label{alg:facility-cubby-slot}
  \begin{algorithmic}[1]
    \State 输入查询键 \(x\)
    \State 计算 \(g_x\gets \pi(x)\)
    \State  \(g_x\) 低位拆分出 Facility 编号 \(r\) 和 Router 桶号 \(b\)
    \State 在表中 facilities[r].D[b] 调用 locate(x)
    \If{Router 返回候选位置 \((C,i)\)}
      \State 读取 C.slots[i] 和 C.meta[i]
      \State 重建候选槽位对应的 \(g_x'\)
      \State 若 \(g_x^{'}=g_x\)，返回命中
    \EndIf
    \State 返回未命中
  \end{algorithmic}
\end{algorithm}

\section{参数产生与地址拆分}

参数包括 \(n\)、k-kick 参数 \(k\)、负载因子 load\_factor 和 Feistel 哈希种子\footnote{Feistel 网络（将输入分为左右两半、多轮轮函数混合的可逆构造）的随机种子，决定可逆置换 \(\pi(x)\) 的具体形式。}。初始化时，代码首先将表规模设置为
\[
 N=\operatorname{next\_pow2}(\max\{2,n\})\footnote{将整数向上取为不小于它的最小 2 的幂，用作表规模 \(N\)。},
\]
因此 \(N\) 始终为 2 的幂，且当 \(n>1\) 时满足 \(N/2<n\leq N\)。随后 choose\_K\footnote{按 \(N\) 计算 Router 桶数 \(K\)，近似为 \((\log_2 N)^2\) 并限制在 \([64,4096]\) 内，再向下取 2 的幂且满足 \(K\le N\)。}根据 \(N\) 生成 \(K\)：先取近似 \((\log_2 N)^2\) 的规模，再限制在 \([64,4096]\) 和 \(K\leq N\) 的范围内，最后向下取 2 的幂\footnote{理论结构要求 \(K\) 为多对数规模，即 \(K=O(\operatorname{polylog}N)\)。当前实现额外要求 \(K\) 为 2 的幂，以便地址拆分与位运算保持一致。}。Facility 数量为 \(\frac{N}{K}\)。

可逆置换由三轮 Feistel 实现, 对任意键 \(x\)，系统先计算 \(g_x=\pi(x)\)，再取其低 \(\log N\) 位：
\[
 g = g_x \mathbin{\mathrm{and}} (N-1).
\]
由于 \(N\) 是 2 的幂，上式等价于 \(g_x\bmod N\)。随后按照
\[
 r=\lfloor g/K\rfloor,\qquad b=g\bmod K.
\]
这样获得 Facility 编号\(r\) 和 Router 桶号 \(b\), 如算法~\ref{alg:route-split} 所示，当 \(N\)、\(K\)、Facility 编号或者 Cubby 容量发生改变时，既有的 payload 与 MetaEntry 的恢复上下文就会随之改变，所以扩缩容迁移不能简单地复制旧槽位的内容。

\begin{algorithm}[htbp]
 \caption{键映射到 Facility 及 Router 桶的流程}
 \label{alg:route-split}
 \begin{algorithmic}[1]
  \State 输入键 \(x\)
  \State \(g_x\gets \pi(x)\)
  \State \(g\gets g_x\mathbin{\mathrm{and}}(N-1)\)
  \State \(r\gets \lfloor g/K\rfloor\)
  \State \(b\gets g\bmod K\)
  \State 返回第 \(r\) 个 Facility 及其中的 Router D[b]
 \end{algorithmic}
\end{algorithm}

\section{Facility、Cubby 与空闲槽结构}

Facility 包括四部分：按 tiered 分组的 Cubbies 集合、tail 及其所有权信息、tiered 上界参数、长度为 \(K\) 的 Router 数组。每个 Router 桶维护“键到候选槽位”的局部映射；每个 Facility 在任一时刻只保留一个 tail，这样就能保证系统装填因子能到达 1。

Cubby 包括tiered(是第几层) 、容量、当前元素数、槽位数组、占用槽索引、空闲槽检索结构和变长元数据数组。槽位只保存商压缩后的 payload，不直接保存完整键；完整恢复依赖 MetaEntry、槽位下标与 Facility 上下文。占用槽索引用于迁移扫描和删除回填，避免在空槽上做无效遍历。

FreeSlotTree\footnote{Cubby 内基于线段树维护各区间空闲槽计数的结构，用于快速查找空槽并更新占用状态。}用线性化的线段树去记录 Cubby 内各区间空闲槽数量，初始化的时候，叶层宽度选取不低于 Cubby 容量的最小 2 的幂，合法槽位的叶子计数设为 1。调用 find\_free\footnote{在 Cubby 的空闲槽线段树中自根向下查找，返回一个可用槽位下标。}时，系统从根节点往下找还有空闲计数的子区间，最后返回一个可用槽位；调用 mark\_used\footnote{将指定槽位标记为已占用，并自底向上更新线段树各区间空闲计数。}或者 mark\_free\footnote{将指定槽位标记为空闲，并更新线段树计数。}后面就顺着父节点去更新计数，这个结构没有做到理论上的低深度压缩位图效果，而是用比较直观的线段树形式来做守护工作，这样就能在插入，删除以及重建的时候检查空闲状况。

Facility中存在两种不同的层次概念，需要加以区分，其一为 k-kick\footnote{Cubby 内部的层级踢出规则，将顺序槽位抽象为 \(k+1\) 层区间并按优先级驱逐。}，这是在一个cubby内部的分层树状管理；其二为 tiered\footnote{Facility 层面的 Cubby 分级组织，按层级 \(j\) 管理不同容量与数量的 Cubby 以支持扩缩容。}，这是用来管理扩缩容实现的分层次cubbies,是论文理论的混淆点与核心难点。

\subsection{tail 管理}

tail 是各 Facility 中唯一允许未填满的 Cubby，在执行插入操作之前，系统会通过调用 ensure\_tail\footnote{保证当前 Facility 存在未满的 tail Cubby；若当前 tail 已满则将其挂入对应 tier 并新建 tier 0 的 tail。}来确保存在一个未满的 tail。如果当前的 tail 已经填满，则代码会把它移至 $tiers[tail\_tier]$ 位置，并且创建一个新的 tier 0 的 tail，其大致步骤如算法~\ref{alg:ensure-tail} 所示：

\begin{algorithm}[htbp]
  \caption{tail Cubby 维护}
  \label{alg:ensure-tail}
  \begin{algorithmic}[1]
    \State  Facility \(F\)
    \If{\(F.tail\) 存在且未满}
      \State 返回
    \EndIf
    \If{\(F.tail\_owned\) 存在且已满}
      \State 将旧 tail 挂入 F.tiers[F.tail\_tier]
      \State 清空 F.tail
    \EndIf
    \State 创建 1-tiered Cubby，初始化 slots、meta、free\_slots
    \State 令 F.tail 指向该 Cubby，并由 F.tail\_owned 持有
  \end{algorithmic}
\end{algorithm}

删除操作时，如果删除的元素不在 tail Cubby 中，删除后会从tail Cubby 中随机取一个元素插入对应的Cubby中，并同步更新 payload、MetaEntry、occupied、free\_slots 和 Router。如果 tail Cubby 中没有元素，将获取一个同样 tiered 层的 cubby 设为 tail（可能触发更高 tiered 层 cubby 的 rebuild\_down）。这样能始终维持非tail cubby 满载这个不变量。

\subsection{j-tiered cubbies 数量与重建调度}

本系统代码实现了 Facility 中多 tiered Cubbies 的基础组织、数量检测和拆分/合并函数。目标数量由 target\_tier\_count\footnote{按层级 \(j\) 与当前元素规模计算该层 Cubby 的目标个数 \(t_j\)。}(j) 计算：
\[
  t_j=\left\lfloor
  \frac{(\log^{(j)} n)^2}{(\log^{(j+1)} n)^2}
  \right\rfloor,
\]
并至少取 1。Cubby 容量由 cubby\_capacity\footnote{由 \(K\)、\(N\) 与 tier 计算该层 Cubby 的槽位容量，并限制在 \([4,K]\) 内。}(K,N,tier) 派生：
\[
  |r_j|\approx \frac{K}{(\log^{(j)}N)^2},
\]
同时限制在 \([4,K]\) 范围内。

插入成功后，maybe\_schedule\_rebuild\footnote{检查各 tier 的 Cubby 数量是否偏离 \(t_j\)，必要时将合并或拆分任务加入后台重建队列。}检查各 tiered 的 Cubby 数量。若 j-tiered cubbies 的数量大于 \(3t_j\)，调用 rebuild\_down\footnote{取出 \(t_j\) 个 \(j\) 层 Cubby，恢复其中键后合并为一个 \((j+1)\) 层 Cubby 并更新 Router。}，将 \(t_j\) 个 j-tiered cubbies 合并为一个 \((j+1)\)-tiered cubby（数据需要恢复，然后重新依次插入）；若数量小于 \(\frac{t_j}{2}\)，且 \((j+1)\)-tiered cubbies 非空，则调度 rebuild\_up\footnote{取出一个 \((j+1)\) 层 Cubby，恢复键后拆成 \(t_j\) 个 \(j\) 层 Cubby 并重新插入。}，将其中一个 \((j+1)\)-tiered cubby 拆分为 \(t_j\) 个 j-tiered cubbies。

\begin{algorithm}[htbp]
  \caption{Facility tier 重建调度}
  \label{alg:tier-schedule}
  \begin{algorithmic}[1]
    \For{\(j=0,1,\ldots,F.max\_tier-1\)}
      \State \(t_j\gets \Call{target\_tier\_count}{j}\)
      \State \(cnt\gets |F.tiers[j]|\)
      \If{\(cnt>3t_j\)}
        \State 调度 \Call{rebuild\_down}{ $F,j$ }
      \ElsIf{\(t_j\ge 2\) 且 \(cnt<t_j/2\) 且 tier \(j+1\) 非空}
        \State 调度 \Call{rebuild\_up}{$F,j$}
      \EndIf
    \EndFor
  \end{algorithmic}
\end{algorithm}

rebuild\_down 会从 j-tiered cubby 中取出 \(t_j\) 个 Cubby，扫描其占用槽并恢复键，再创建一个 \((j+1)\)-tiered Cubby 重新写入这些键；rebuild\_up 则从 \((j+1)\)-tiered 层取出一个 Cubby，恢复其中所有键，创建 \(t_j\) 个 j-tiered Cubby 并将恢复的原始键全部插入其中。两者都会更新 Router，使 Router 记录指向新的 Cubby 和 Slot。当前实现具备拆分/合并的数据路径，但其调度粒度和容量控制仍属于工程化实现，不能直接推出理论均摊界。

\section{MiniArray 与 MetaEntry}

\subsection{MiniArray 变长元数据存储结构}
MiniArray 是 Cubby 组件中用于存储各槽位变长元数据的核心结构，通过多组向量协同维护元数据的长度、偏移、原始内容与连续打包存储区域，实现变长数据的紧凑管理。

MiniArray 各成员变量的功能定义如下：
\begin{itemize}
    \item bitlens\_[i]：记录第 $i$ 个元数据条目的二进制位长度；
    \item offsets\_[i]：标记第 $i$ 个元数据条目在连续位缓冲区中的起始比特位置；
    \item entries\_[i]：保存元数据条目的未压缩原始副本；
    \item buf\_：存储所有元数据经重新打包后的连续比特缓冲区。
\end{itemize}


\subsubsection{读取流程}

\begin{algorithm}
    \caption*{MiniArray 读取流程}
    \begin{algorithmic}[1]
        \Function{access}{$i$}
        \State $bl \gets bitlens\_[i]$
        \State $start \gets offsets\_[i]$
        \State 从 $buf\_$ 中起始位置 $start$ 取出长度为 $bl$ 的比特数据
        \State 将数据组织为 64 位字数组并返回
        \EndFunction
    \end{algorithmic}
\end{algorithm}
读取操作由 access\footnote{按槽位下标从连续缓冲区读出该条变长元数据的比特串。}完成，基于位长度与起始偏移量提取指定元数据并返回标准格式结果：

\subsubsection{更新流程}
MiniArray 的更新操作由 update\footnote{修改指定槽位的元数据并调用 repack 重算偏移与缓冲区。}完成，针对指定槽位完成元数据修改，并触发全局重打包以保证存储布局一致性，流程如下：

\begin{algorithm}
    \caption*{MiniArray 更新流程}
    \begin{algorithmic}[1]
        \Procedure{update}{$i$, $bits$, $bitlen$}
            \If{$i$ 越界} \State 报错 \EndIf
            \State $bitlens\_[i] \gets bitlen$
            \State $entries\_[i] \gets bits$
            \State \Call{repack}{}
        \EndProcedure
        \Procedure{repack}{} \Comment{按位长重算偏移并写回 buf\_}
            \State 根据 bitlens\_ 重新计算所有条目的偏移量 offsets\_
            \State 为连续缓冲区 buf\_ 分配内存空间
            \State 将所有元数据条目按偏移逐位写入 buf\_
        \EndProcedure
    \end{algorithmic}
\end{algorithm}

在理论设计中，MiniArray 可通过内存 B 树与查表法实现均摊常数时间的更新操作\cite{bender2022otsh}；打包式变长存储在 PaCHash\cite{kurpicz2023pachash} 等工作中亦有类似思路。本实现保留变长条目管理、连续比特打包与位数统计的接口语义，以全局重打包替代复杂动态结构，在保证正确性的前提下简化逻辑。

\subsection{MetaEntry 元数据结构}
MetaEntry\footnote{槽位元数据的结构化字段集合，包含 delta、mid\_bits、router\_bits、insert\_bits 等，配合 payload 恢复完整置换值 \(g_x\)。}是 Cubby 槽位元数据的结构化表示，用于存储哈希定位、路由校验、插入追踪所需的核心信息，MetaEntry各字段含义如下：
\begin{itemize}
    \item delta：计算公式为 $g_i - \text{slot\_index}$，用于根据实际槽位恢复哈希计算得到的 Cubby 偏好槽位；
    \item mid\_bits：计算公式为 $g_k / |I|$，用于恢复局部桶区间内的中间比特信息；
    \item router\_bits：路由校验指纹字段，由 $\text{splitmix64}(\text{key} \oplus \text{gx\_pi} \oplus \text{常量})$ 计算后截断为 16 位得到；
    \item insert\_bits：插入过程标签，主要用于记录数据插入过程中的 k-kick 搬移步数信息。
\end{itemize}

\subsubsection{位宽配置与编码规则}
MetaCodecLayout\footnote{由 \(K\) 与 Cubby 容量派生 mid\_bits、router\_bits、insert\_bits 等字段的编码位宽。}根据哈希表参数 $K$ 与 Cubby 容量动态确定各字段编码位宽：

\begin{algorithm}
    \caption*{位宽配置}
    \begin{algorithmic}[1]
    \State $mid\_w \gets \text{bit\_width}((K - 1) / \text{cubby\_capacity})$
    \State $rout\_w \gets 16$
    \State $ins\_w \gets 8$
    \end{algorithmic}
\end{algorithm}

\subsubsection{MetaEntry 生成逻辑}
make\_meta\_entry\footnote{根据键、置换值、槽位下标及表参数生成 delta、mid\_bits、router\_bits、insert\_bits 等字段。}是 MetaEntry 结构体的核心生成函数，基于哈希计算结果与插入参数完成字段赋值，核心逻辑如下：

\begin{algorithm}
    \caption*{MetaEntry 构造}
    \begin{algorithmic}[1]
    \State $g\_\text{star} \gets g_x \& (N - 1)$ \hfill // 计算全局哈希种子（取模 N）
    \State $g\_k \gets g\_\text{star} \& (K - 1)$ \hfill // 提取组内索引（取模 K）
    \State $g\_i \gets g_k \bmod cubby\_capacity$ \hfill // 计算在 cubby 中的偏移
    \State $delta \gets g_i - slot\_index$ \hfill // 计算与当前槽位的偏移差
    \State $mid\_bits \gets g_k / cubby\_capacity$ \hfill // 提取中间位（组编号）
    \State $router\_bits \gets \text{截断}(\text{hash}(key, g_x),\ 16)$ \hfill // 路由信息（16 位哈希截断）
    \State $insert\_bits \gets ins\_tag \& \text{0xff}$ \hfill // 插入标签（低 8 位）
    \end{algorithmic}
\end{algorithm}

函数参数 facility\_r 仅作为上下文参数传入，不直接参与字段计算，该参数将在数据恢复阶段与元数据配合完成哈希值的完整重建。

\section{Router 组件}

Router要维护键到 \((\text{Cubby}^*, \textit{slot})\) 的映射关系，其内部条目包含原始键，目标位置以及由 key\_bits\footnote{为键生成 Router 前缀树所用的 64 位随机比特串。}得到的随机比特串，而节点则保存左右孩子，叶子条目的下标，碰撞桶的下标以及分裂位 split\_bit。

Router并非按键值大小来组织结构，而是凭借随机比特串形成 Patricia 风格的二叉前缀树\footnote{Patricia 树是一种压缩前缀树：内部结点仅在有分支处按某一比特位分裂，叶结点存完整键或条目。}，此前缀树仅在Router内部用以识别随机比特串，执行查询时，系统从根节点出发，遵照split\_bit获取查询键中的随机比特，并决定向左或向右子树探寻；直至抵达叶节点才对比原本的键；要是存在多个键其64位随机比特串完全一致，则需执行碰撞桶\footnote{Patricia 树中随机比特串相同的多个键所落入的线性扫描容器，用于消解哈希碰撞。}的线性扫描操作。

\begin{algorithm}[htbp]
  \caption{Router 查询流程}
  \label{alg:router-locate}
  \begin{algorithmic}[1]
    \State 输入键 \(x\)
    \State \(\textit{bits}\gets\Call{key\_bits}{x}\)，\(\textit{steps}\gets0\)，\(u\gets root\)
    \While{\(u\) 为合法结点}
      \State \(\textit{steps}\gets\textit{steps}+1\)
      \If{\(u\) 为碰撞桶}
        \State 在桶内逐项比较原始键
        \State 返回命中位置或未命中，以及 \(\textit{steps}\)
      \ElsIf{\(u\) 为叶子}
        \State 比较叶子 entry 中的原始键
        \State 返回命中位置或未命中，以及 \(\textit{steps}\)
      \Else
        \State 按 \(\textit{bits}\) 在 split\_bit 处的取值进入左/右子树
      \EndIf
    \EndWhile
    \State 返回未命中
  \end{algorithmic}
\end{algorithm}

插入或者删除 Router 记录之后，代码会调用 rebuild\footnote{重新计算各 entry 的随机比特串，重建 Patricia 前缀树并更新编码位长统计。}，并利用 BitWriter 遍历结构来估算 encoded\_bits\_。所以 bits\_total\footnote{返回 Router 前缀树逻辑编码的累计比特数，用于空间统计。}可以用来了解路由结构的逻辑编码规模，但它并不能表示真实的序列化存储或者运行时的内存占用情况，按照理论中的 Local Query Router 理应可以在紧凑 bitstream 上做到常数时间的动态更新；不过当下的实现则是用整体重建的方式来取代这种复杂的结构。

\section{商压缩与键恢复}

商压缩由槽内 payload、MiniArray 中的 MetaEntry 以及当前表上下文共同完成。设待处理键为 \(x\)，\(g_x=\pi(x)\)，\(\ell=\log N\)。槽位 slots[i] 保存
\[
  payload=g_x \gg \ell,
\]
即去掉低 \(\ell\) 位后的高位商。低 \(\ell\) 位参与地址拆分和局部恢复，不直接完整保存。

令
\[
  g^\ast=g_x\mathbin{\mathrm{and}}(N-1),\quad
  r=\lfloor g^\ast/K\rfloor,\quad
  g_k=g^\ast\mathbin{\mathrm{and}}(K-1).
\]
给定 Cubby 容量 \(|I|\) 和实际槽位 \(i\)，代码计算
\[
  g_i=g_k\bmod |I|,\quad
  delta=g_i-i,\quad
  mid\_bits=\lfloor g_k/|I|\rfloor .
\]
恢复时再按
\[
  g_i=i+delta,\quad
  g_k=mid\_bits\times |I|+g_i,\quad
  g^\ast=rK+g_k,\quad
  g_x=(payload\ll \ell)\,|\,g^\ast
\]
重建完整置换值。若需要恢复原始键，则继续调用 Feistel 逆置换 \(\pi^{-1}\)。

\begin{figure}[htbp]
  \centering
  \setlength{\tabcolsep}{6pt}
  \begin{tabular}{|c|c|c|c|}
    \hline
    rest / payload & Facility 编号 \(r\) & 中间位 mid\_bits & 偏好槽位 \(g_i\) \\
    \hline
    \(g_x >> \log N\) & \(\lfloor g^\ast/K\rfloor\) & \(\lfloor g_k/|I|\rfloor\) & \(g_k\bmod |I|\) \\
    \hline
    存入 slots[i] & 由地址上下文提供 & 存入 MetaEntry & 由 delta 与槽位恢复 \\
    \hline
  \end{tabular}
  \caption{\(\pi(x)\) 在商压缩实现中的位段用途}
  \label{fig:gx-bit-layout}
\end{figure}

payload 以及 MetaEntry 都依靠表规模，局部桶数，Facility 编号和 Cubby 容量，执行扩缩容迁移之后，元素需先还原成原本的键，然后依照新表的状态加以重新写入，不可直接拷贝旧槽位里的 payload 和 meta。

\begin{algorithm}[htbp]
  \caption{从槽位恢复原始键}
  \label{alg:recover-key}
  \begin{algorithmic}[1]
    \State 输入 Facility 编号 \(r\)、Cubby \(C\)、槽位 \(i\)
    \If{\(i\) 越界或 C.slots[i]} 为空
      \State 返回空
    \EndIf
    \State 从 C.meta[i] 解码 MetaEntry
    \State 根据 \(r,N,K,C.capacity,i,\textit{payload}\) 重建 \(g_x\)
    \If{重建失败}
      \State 返回空
    \EndIf
    \State 返回 \(\pi^{-1}(g_x)\)
  \end{algorithmic}
\end{algorithm}


\section{双表渐进迁移}

论文提及了两种程度的扩容和缩容，其中一种是多层次的 j-tiered cubbies 设计，这已经实现，另一种是在大规模的扩容缩容，完全是新的表与旧的表进行迁移，例如扩容时，会将旧表中的两个Facility的数据插入新表的一个Facility中，同理缩容相同的逻辑。当系统哈希表正在进行迁移时，会对应旧表和新表依次执行对应的插入、删除、查询操作，如果操作完成就提前终止，想要实现这样的无感迁移，就需要实现渐进式的双表迁移\footnote{当前系统并没有完整地实现渐进式双表迁移，只实现了代码框架，并没有数据支撑；渐进式扩缩容的思想与动态可调整哈希表\cite{bender2024sosa}等近年方法一脉相承。}，具体流程如算法~\ref{alg:migrate} 所示。

\begin{algorithm}[htbp]
    \caption{单次迁移预算}
    \label{alg:migrate}
    \begin{algorithmic}[1]
        \If{old\_} 不存在
        \State 返回
        \EndIf
        \State 取 migrate\_progress\_ 指向的旧 Facility
        \State 扫描该 Facility 的 tail 和所有 tier Cubby，恢复占用键集合 \(S\)
        \For{\(x\in S\)}
        \State 按新表参数计算 \(g_x=\pi(x)\)，定位 active Facility
        \If{\(x\) 已存在于 active 表}
        \State 跳过
        \Else
        \State 调用 ensure\_tail 和 kkick\_insert\footnote{在指定 Cubby 内执行有界 k-kick 插入，返回是否成功及搬移次数。}重新写入
        \EndIf
        \EndFor
        \State 清空该旧 Facility，推进迁移进度
        \State 若所有旧 Facility 均已处理，则释放 old\_
    \end{algorithmic}
\end{algorithm}

系统经由 maybe\_resize\_locked\footnote{在持锁状态下根据负载因子判断是否需要扩容或缩容。}来判断是否执行扩缩容操作，当下不存在 old 表，而且元素数量超出 \(load\_factor\times N\) 时，把目标规模设置成 \(2N\)。如果元素数量不为零，又小于 \(N/4\) 的话，则把目标规模设定为 \(\max\{2,N/2\}\)，一旦发生这种情况，start\_migration\footnote{将当前 active 表移入 old\_，按新规模重建 active 表并初始化迁移进度。}就会清空重建调度队列，并把 active 表移至 old\_，然后按照新的 \(N\) 和 \(K\) 创建一个新的 active 表，这个过程中，元素总数 n\_ 保持不变，新增的数据只会被添加到 active 表当中。

迁移由 run\_migrate\_budget\footnote{每次操作后最多迁移一个旧 Facility：恢复其中键并插入 active 表，推进 migrate\_progress\_。}推动，每当插入或者删除操作结束之后，系统最多只会处理一个旧的 Facility。在处理过程中不依靠旧 Router，而是要遍历旧 Facility 的 tail 及其所有的 tier Cubby，经由 recover\_key\_from\_slot\footnote{从槽位 payload 与 MetaEntry 重建置换值并求逆，得到原始键。}（算法~\ref{alg:recover-key}）来重建原本的键，然后按照新的表参数重新确定位置并执行插入到 active 表的操作，如果某个键已在 active 表中有对应项，则直接忽略这个键，当一个旧 Facility 的处理完毕之后，它的 Router，tiers 和 tail 都会被清除。

该迁移方式避免扩缩容时一次性重哈希全部元素。其限制是：查询不会消耗迁移预算；若扩缩容后长期没有插入或删除，old 表会继续保留。此外，由于 payload 和 MetaEntry 依赖新表上下文，迁移必须恢复键并重新插入。

\section{插入路径与 k-kick 实现}

外层函数加锁后调用 insert\_no\_lock\footnote{在无额外加锁前提下完成插入、扩缩容判断及后台重建/迁移预算。}，成功插入后检查扩缩容，随后执行一个后台重建预算和一个迁移预算。真正写入发生在 kkick\_insert 中。

插入流程如下：先算出 \(g_x = \pi(x)\)，找到 active 表里对应的 Facility 和 Router 桶；如果 active 表已有此键，则直接返回 inserted=false。若是处于迁移窗口，则还要查看 old 表，防止同一个键在迁移过程中被重复写入，确定键不存在之后，系统会调用 ensure\_tail 来准备好 tail Cubby，并在其中执行 k-kick 操作。

当前 k-kick 是单 Cubby 内的有界踢出链。探测位置由
\[
  \operatorname{probe\_slot}\footnote{按当前键、Cubby 容量与踢出步数 \(s\) 计算本步探针槽位。}(x,|I|,s)
  =
  \operatorname{splitmix64}(\pi(x)\oplus s\cdot c)\bmod |I|
\]
得到，其中 \(s\) 为当前步数，\(c\) 为固定常量。若探测槽为空，系统写入 payload 和 MetaEntry，更新 occupied、free\_slots 和 Router；若探测槽已占用，则先恢复被踢出的旧键，删除旧 Router 记录，再用当前键覆盖该槽，并将旧键作为下一轮待插入键。经过 \(k\) 轮仍未结束时，系统通过 FreeSlotTree 选择空槽兜底。

\begin{algorithm}[htbp]
  \caption{tail Cubby 内的有界踢出插入}
  \label{alg:kkick}
  \begin{algorithmic}[1]
    \State 输入 Facility \(F\)、tail Cubby \(C\)、待插入键 \(x\)
    \State \(\textit{cur}\gets x\)
    \For{\(s=0,1,\ldots,k\)}
      \State \(pos\gets\Call{probe\_slot}{\textit{cur},C.capacity,s}\)
      \If{C.slots[pos] 为空}
        \State 写入 \(\textit{cur}\) 的 payload 和 MetaEntry
        \State 更新 occupied、free\_slots 与 Router
        \State 返回成功，移动次数为 \(s\)
      \EndIf
      \State 从 \(pos\) 恢复被踢出的键 \(\textit{kicked}\)
      \If{恢复失败}
        \State 返回失败
      \EndIf
      \State 删除 \(\textit{kicked}\) 的旧 Router 记录
      \State 用 \(\textit{cur}\) 覆盖 \(pos\)，并登记新 Router 记录
      \State \(\textit{cur}\gets \textit{kicked}\)
    \EndFor
    \State 通过 FreeSlotTree 查找空槽，写入 \(\textit{cur}\)
    \State 返回成功，移动次数为 \(k+1\)
  \end{algorithmic}
\end{algorithm}

理论 k-kick 树需要维护更完整的深度、饱和状态和合法驱逐对象。当前实现保留了“单 Cubby 内有限交换”的核心接口，但用伪随机探测链和空槽兜底替代理论结构。因此，实验中的 kick\_count 表示当前代码中的移动次数，不能直接等同于理论交换成本。

\section{查询与删除路径}

查询路径仅读取结构，不会推进迁移预算，系统首先在 active 表里定位 Router 桶，经由 Router::locate\footnote{在 Router 前缀树中查找键，返回候选 \((\text{Cubby},\text{slot})\) 及访问步数。}得到候选位置，然后用 slot\_pi\_matches\footnote{从候选槽位解码元数据并重建 \(\pi(x)\)，与查询键的置换值比对以判定命中。}解码 meta，重新创建 \(\pi(x)\) 并做比较，如果 active 表未命中而且 old 表存在，就对 old 表重复前面的过程，Router 访问步数会被记入 router\_probe\_steps。

删除路径先执行与查询相同的定位和槽位校验。若命中 active 表中的槽位，系统清空 slots[slot]，将对应 MiniArray 条目更新为空，移除 occupied 中的槽位下标，调用 mark\_free 标记空闲，并从 Router 中删除该键。如果被删除的 Cubby 不是 tail，系统会从 tail 的 occupied 中取出一个元素，恢复其原始键后搬移到刚刚空出的槽位，并重新写入 payload、MetaEntry 和 Router 记录。

\begin{algorithm}[htbp]
  \caption{查询命中判定}
  \label{alg:query}
  \begin{algorithmic}[1]
    \State 输入查询键 \(x\)
    \State 在 active 表中定位 \((r,b)\)，调用 D[b].locate(x)
    \If{Router 返回候选槽位 \((C,i)\)}
      \State 调用 slot\_pi\_matches 校验候选槽位
      \State 若校验通过，返回命中
    \EndIf
    \If{active 表未命中且 old 表存在}
      \State 在 old 表中重复定位与槽位校验
    \EndIf
    \State 返回最终查询结果
  \end{algorithmic}
\end{algorithm}

当处在迁移的状态时，插入只写入 active 表，但查询和删除必须同时考虑 active 表和 old 表。删除时若 old 表中仍存在该键，也需要删除，否则后续迁移可能把旧键重新搬入 active 表。
